\title{2025 年高考数学 天津卷}
\date{}
\maketitle

\section*{一、单项选择题}
在每小题给出的四个选项中，只有一项是符合题目要求的

\begin{question}[points=5]
已知集合$U = \{1,2,3,4,5\}$, $A = \{1,3\}$, $B = \{2,3,5\}$, 则$\complement_U (A \cup B) = $（   ）
\begin{choices}
  \item $\{1,2,3,4\}$
  \item $\{2,3,4\}$
  \item $\{2,4\}$
  \item $\{4\}$
\end{choices}
\begin{solution}
\textbf{答案：}D

\textbf{分析：}由集合的并集、补集的运算即可求解.

\textbf{详解：}由$A = \{1,3\}$, $B = \{2,3,5\}$, 则$A \cup B = \{1,2,3,5\}$, 集合$U = \{1,2,3,4,5\}$, 故$\complement_U (A \cup B) = \{4\}$. 故选D.
\end{solution}
\end{question}

\begin{question}[points=5]
设$x \in \mathbb{R}$, 则“$x = 0$”是“$\sin 2x = 0$”的（   ）
\begin{choices}
  \item 充分不必要条件
  \item 必要不充分条件
  \item 充要条件
  \item 既不充分也不必要条件
\end{choices}
\begin{solution}
\textbf{答案：}A

\textbf{分析：}通过判断是否能相互推出，由充分条件与必要条件的定义可得.

\textbf{详解：}由$x = 0 \Rightarrow \sin 2x = \sin 0 = 0$, 则“$x = 0$”是“$\sin 2x = 0$”的充分条件；又当$x = \pi$时, $\sin 2x = \sin 2\pi = 0$, 可知$\sin 2x = 0 \not\Rightarrow x = 0$, 故“$x = 0$”不是“$\sin 2x = 0$”的必要条件. 综上可知, “$x = 0$”是“$\sin 2x = 0$”的充分不必要条件. 故选A.
\end{solution}
\end{question}

\begin{question}[points=5]
已知函数$y = f(x)$的图象如下, 则$f(x)$的解析式可能为（   ）
\begin{choices}
  \item $f(x) = \dfrac{x}{1 - |x|}$
  \item $f(x) = \dfrac{x}{|x| - 1}$
  \item $f(x) = \dfrac{|x|}{1 - x^2}$
  \item $f(x) = \dfrac{|x|}{x^2 - 1}$
\end{choices}
\begin{solution}
\textbf{答案：}D

\textbf{分析：}先由函数奇偶性排除AB，再由$x \in (0,1)$时函数值正负情况可得解.

\textbf{详解：}由图可知函数为偶函数，而函数$f(x) = \dfrac{x}{1-x}$和函数$f(x) = \dfrac{x}{x-1}$为奇函数，故排除选项AB；又当$x \in (0,1)$时$1 - x^2 > 0$, $x^2 - 1 < 0$, 此时$f(x) = \dfrac{x}{1-x^2} > 0$, $f(x) = \dfrac{x}{x^2-1} < 0$, 由图可知当$x \in (0,1)$时, $f(x) < 0$, 故C不符合，D符合. 故选D.
\end{solution}
\end{question}

\begin{question}[points=5]
若$m$为直线，$\alpha, \beta$为两个平面，则下列结论中正确的是（   ）
\begin{choices}
  \item 若$m // \alpha$, $n \subset \alpha$, 则$m // n$
  \item 若$m \perp \alpha$, $m \perp \beta$, 则$\alpha \perp \beta$
  \item 若$m // \alpha$, $m \perp \beta$, 则$\alpha \perp \beta$
  \item 若$m \subset \alpha$, $\alpha \perp \beta$, 则$m \perp \beta$
\end{choices}
\begin{solution}
\textbf{答案：}C

\textbf{分析：}根据线面平行的定义可判断A的正误，根据空间中垂直关系的转化可判断BCD的正误.

\textbf{详解：}对于A，若$m // \alpha$, $n \subset \alpha$, 则$m, n$可平行或异面，故A错误；对于B，若$m \perp \alpha$, $m \perp \beta$, 则$\alpha // \beta$, 故B错误；对于C，两条平行线有一条垂直于一个平面，则另一个必定垂直这个平面，现$m // a$, $m \perp \beta$, 故$a \perp \beta$, 故C正确；对于D，$m \subset \alpha$, $\alpha \perp \beta$, 则$m$与$\beta$可平行或相交或$m \subset \beta$, 故D错误. 故选C.
\end{solution}
\end{question}

\begin{question}[points=5]
下列说法中错误的是（   ）
\begin{choices}
  \item 若$X \sim N(\mu, \sigma^2)$, 则$P(X \le \mu - \sigma) = P(X \ge \mu + \sigma)$
  \item 若$X \sim N(1, 2^2)$, $Y \sim N(2, 2^2)$, 则$P(X < 1) < P(Y < 2)$
  \item $r$越接近1，相关性越强
  \item $r$越接近0，相关性越弱
\end{choices}
\begin{solution}
\textbf{答案：}B

\textbf{分析：}根据正态分布以及相关系数的概念直接判断即可.

\textbf{详解：}对于A，根据正态分布对称性可知，$P(X \le \mu - \sigma) = P(X \ge \mu + \sigma)$, A说法正确；对于B，根据正态分布对称性可知，$P(X < 1) = P(Y < 2) = 0.5$, B说法错误；对于C和D，相关系数$r$越接近0，相关性越弱，越接近1，相关性越强，故C和D说法正确. 故选B.
\end{solution}
\end{question}

\begin{question}[points=5]
$S_n = -n^2 + 8n$, 则数列$\{a_n\}$的前12项和为（   ）
\begin{choices}
  \item 112
  \item 48
  \item 80
  \item 64
\end{choices}
\begin{solution}
\textbf{答案：}C

\textbf{分析：}先由题设结合$a_n = S_n - S_{n-1}$求出数列$\{a_n\}$的通项公式，再结合数列$\{a_n\}$各项正负情况即可求解.

\textbf{详解：}因为$S_n = -n^2 + 8n$, 所以当$n=1$时, $a_1 = S_1 = -1 + 8 \times 1 = 7$; 当$n \ge 2$时, $a_n = S_n - S_{n-1} = (-n^2 + 8n) - [-(n-1)^2 + 8(n-1)] = -2n + 9$. 经检验$a_1 = 7$满足上式, 所以$a_n = -2n + 9$ ($n \in \mathbb{N}^*$). 令$a_n = -2n + 9 \ge 0 \Rightarrow n \le 4$, $a_n = -2n + 9 \le 0 \Rightarrow n \ge 5$. 设数列$\{|a_n|\}$的前$n$项和为$T_n$, 则数列$\{|a_n|\}$的前4项和为$T_4 = S_4 = -4^2 + 8 \times 4 = 16$. 数列$\{|a_n|\}$的前12项和为$T_{12} = a_1 + a_2 + \cdots + a_{12} = a_1 + a_2 + a_3 + a_4 - a_5 - a_6 - \cdots - a_{12} = 2S_4 - S_{12} = 2 \times 16 - (-12^2 + 8 \times 12) = 80$. 故选C.
\end{solution}
\end{question}

\begin{question}[points=5]
函数$f(x) = 0.3^x - \sqrt{x}$的零点所在区间是（   ）
\begin{choices}
  \item $(0, 0.3)$
  \item $(0.3, 0.5)$
  \item $(0.5, 1)$
  \item $(1, 2)$
\end{choices}
\begin{solution}
\textbf{答案：}B

\textbf{分析：}利用指数函数与幂函数的单调性结合零点存在性定理计算即可.

\textbf{详解：}由指数函数、幂函数的单调性可知：$y = 0.3^x$在$\mathbb{R}$上单调递减, $y = \sqrt{x}$在$[0, +\infty)$单调递增, 所以$f(x) = 0.3^x - \sqrt{x}$在定义域上单调递减. 显然$f(0) = 1 > 0$, $f(0.3) = 0.3^{0.3} - \sqrt{0.3} > 0$, $f(0.5) = 0.3^{0.5} - \sqrt{0.5} < 0$, 所以根据零点存在性定理可知$f(x)$的零点位于$(0.3, 0.5)$. 故选B.
\end{solution}
\end{question}

\begin{question}[points=5]
$f(x) = \sin(\omega x + \varphi)$ ($\omega > 0$, $-\pi < \varphi < \pi$), 在$\left[-\dfrac{5\pi}{12}, \dfrac{\pi}{12}\right]$上单调递增, 且$x = \dfrac{\pi}{12}$为它的一条对称轴, $\left(\dfrac{\pi}{3}, 0\right)$是它的一个对称中心, 当$x \in \left[0, \dfrac{\pi}{2}\right]$时, $f(x)$的最小值为（   ）
\begin{choices}
  \item $-\dfrac{\sqrt{3}}{2}$
  \item $-\dfrac{1}{2}$
  \item 1
  \item 0
\end{choices}
\begin{solution}
\textbf{答案：}A

\textbf{分析：}利用正弦函数的对称性得出$\omega = 4n + 2$, 根据单调性得出$0 < \omega \le 2$, 从而确定$\omega$, 结合对称轴与对称中心再求出$\varphi$, 得出函数解析式, 利用整体思想及正弦函数的性质即可得解.

\textbf{详解：}设$f(x)$的最小正周期为$T$, 根据题意有$\begin{cases} \dfrac{\pi}{12}\omega + \varphi = \dfrac{\pi}{2} + 2k\pi \\ \dfrac{\pi}{3}\omega + \varphi = m\pi \end{cases}$, $(m, k \in \mathbb{Z})$. 由正弦函数的对称性可知$\dfrac{\pi}{3} - \dfrac{\pi}{12} = \dfrac{(2n+1)T}{4}$ ($n \in \mathbb{Z}$), 即$\dfrac{\pi}{4} = \dfrac{2n\pi + \pi}{2\omega}$, $\therefore \omega = 4n + 2$. 又$f(x)$在$\left[-\dfrac{5\pi}{12}, \dfrac{\pi}{12}\right]$上单调递增, 则$\dfrac{T}{2} \ge \dfrac{\pi}{12} - \left(-\dfrac{5\pi}{12}\right)$, $\therefore \dfrac{\pi}{\omega} \ge \dfrac{\pi}{2} \Rightarrow 0 < \omega \le 2$. $\therefore \omega = 2$, 则$\begin{cases} \varphi = \dfrac{\pi}{3} + 2k\pi \\ \varphi = m\pi - \dfrac{2\pi}{3} \end{cases}$. $\because \varphi \in (-\pi, \pi)$, $\therefore k=0$, $m=1$时, $\varphi = \dfrac{\pi}{3}$, $\therefore f(x) = \sin\left(2x + \dfrac{\pi}{3}\right)$. 当$x \in \left[0, \dfrac{\pi}{2}\right]$时, $2x + \dfrac{\pi}{3} \in \left[\dfrac{\pi}{3}, \dfrac{4\pi}{3}\right]$, 由正弦函数的单调性可知$f(x)_{\min} = \sin\dfrac{4\pi}{3} = -\dfrac{\sqrt{3}}{2}$. 故选A.
\end{solution}
\end{question}

\begin{question}[points=5]
双曲线$\dfrac{x^2}{a^2} - \dfrac{y^2}{b^2} = 1$ ($a > 0$, $b > 0$)的左、右焦点分别为$F_1$, $F_2$, 以右焦点$F_2$为焦点的抛物线$y^2 = 2px$ ($p > 0$)与双曲线交于另一象限点为$P$, 若$|PF_1| + |PF_2| = 3|F_1F_2|$, 则双曲线的离心率$e = $（   ）
\begin{choices}
  \item 2
  \item $\sqrt{5}$
  \item $\dfrac{\sqrt{2}+1}{2}$
  \item $\dfrac{\sqrt{5}+1}{2}$
\end{choices}
\begin{solution}
\textbf{答案：}A

\textbf{分析：}利用抛物线与双曲线的定义与性质得出$\begin{cases} |PF_1| = 3c + a \\ |PF_2| = 3c - a = PA \end{cases}$, 根据勾股定理从而确定P的坐标, 利用点在双曲线上构造齐次方程计算即可.

\textbf{详解：}根据题意可设$F_2\left( \dfrac{p}{2}, 0 \right)$, 双曲线的半焦距为$c$, $P(x_0, y_0)$, 则$p = 2c$. 过$F_1$作$x$轴的垂线$l$, 过$P$作$l$的垂线, 垂足为$A$, 显然直线$AF_1$为抛物线的准线, 则$PA = PF_2$. 由双曲线的定义及已知条件可知$\begin{cases} |PF_1| - |PF_2| = 2a \\ |PF_1| + |PF_2| = 6c \end{cases}$, 则$\begin{cases} |PF_1| = 3c + a \\ |PF_2| = 3c - a = PA \end{cases}$. 由勾股定理可知$|AF_1|^2 = y_0^2 = |PF_1|^2 - |PA|^2 = 12ac$. 易知$y_0^2 = 4cx_0$, $\therefore x_0 = 3a$, 即$\dfrac{x_0^2}{a^2} - \dfrac{y_0^2}{b^2} = \dfrac{9a^2}{a^2} - \dfrac{12ac}{c^2 - a^2} = 1$, 整理得$2c^2 - 3ac - 2a^2 = 0 = (2c + a)(c - 2a)$, $\therefore c = 2a$, 即离心率为2. 故选A.
\end{solution}
\end{question}

\section*{二、填空题}
本大题共6个小题，每小题5分，共30分.

\begin{question}[points=5]
已知$i$是虚数单位，则$\left| \dfrac{3+i}{i} \right| = $.
\fillin[{$\sqrt{10}$}]
\begin{solution}
\textbf{答案：}$\sqrt{10}$

\textbf{分析：}先由复数除法运算化简$\dfrac{3+i}{i}$, 再由复数模长公式即可计算求解.

\textbf{详解：}先由题得$\dfrac{3+i}{i} = -i(3+i) = 1 - 3i$, 所以$\left| \dfrac{3+i}{i} \right| = \sqrt{1^2 + (-3)^2} = \sqrt{10}$. 故答案为$\sqrt{10}$.
\end{solution}
\end{question}

\begin{question}[points=5]
在$(x-1)^6$的展开式中, $x^3$项的系数为.
\fillin[{$-20$}]
\begin{solution}
\textbf{答案：}$-20$

\textbf{分析：}根据二项式定理相关知识直接计算即可.

\textbf{详解：}$(x-1)^6$展开式的通项公式为$T_{r+1} = \mathrm{C}_6^r x^{6-r} \cdot (-1)^r$. 当$r=3$时, $T_4 = \mathrm{C}_6^3 x^3 \cdot (-1)^3 = -20x^3$, 即$(x-1)^6$展开式中$x^3$的系数为$-20$. 故答案为$-20$.
\end{solution}
\end{question}

\begin{question}[points=5]
$l_1: x - y + 6 = 0$, 与$x$轴交于点$A$, 与$y$轴交于点$B$, 与$(x+1)^2 + (y-3)^2 = r^2$交于$C$、$D$两点, $|AB| = 3|CD|$, 则$r = $.
\fillin[{$2$}]
\begin{solution}
\textbf{答案：}$2$

\textbf{分析：}先根据两点间距离公式得出$|AB| = 6\sqrt{2}$, 再计算出圆心到直线的距离$d$, 根据弦长公式$|CD| = 2\sqrt{r^2 - d^2}$列等式求解即可.

\textbf{详解：}因为直线$l_1: x - y + 6 = 0$与$x$轴交于$A(-6,0)$, 与$y$轴交于$B(0,6)$, 所以$|AB| = \sqrt{6^2 + 6^2} = 6\sqrt{2}$, 所以$|CD| = 2\sqrt{2}$. 圆$(x+1)^2 + (y-3)^2 = r^2$的半径为$r$, 圆心$(-1,3)$到直线$l_1: x - y + 6 = 0$的距离为$d = \dfrac{|-1 - 3 + 6|}{\sqrt{2}} = \sqrt{2}$, 故$|CD| = 2\sqrt{r^2 - d^2} = 2\sqrt{r^2 - (\sqrt{2})^2} = 2\sqrt{2}$, 解得$r=2$. 故答案为2.
\end{solution}
\end{question}

\begin{question}[points=5]
小桐操场跑圈，一周2次，一次5圈或6圈. 第一次跑5圈或6圈的概率均为0.5，若第一次跑5圈，则第二次跑5圈的概率为0.4，6圈的概率为0.6；若第一次跑6圈，则第二次跑5圈的概率为0.6，4圈的概率为0.4. 小桐一周跑11圈的概率为；若一周至少跑11圈为动量达标，则连续跑4周，记合格周数为$X$, 则期望$E(X) = $.
\fillin[{$0.6；3.2$}]
\begin{solution}
\textbf{答案：}$0.6；3.2$

\textbf{分析：}先根据全概率公式计算求解空一，再求出概率根据二项分布数学期望公式计算求解.

\textbf{详解：}设小桐一周跑11圈为事件$A$, 设第一次跑5圈为事件$B$, 设第二次跑5圈为事件$C$, 则$P(A) = P(B)P(C|B) + P(\overline{B})P(\overline{C}|\overline{B}) = 0.5 \times 0.6 + 0.5 \times 0.6 = 0.6$；若至少跑11圈为运动量达标为事件$D$, $P(D) = P(A) + P(B)P(\overline{C}|B) = 0.6 + 0.5 \times 0.4 = 0.8$, 所以$X \sim B(4, 0.8)$, $E(X) = 4 \times 0.8 = 3.2$. 故答案为0.6；3.2.
\end{solution}
\end{question}

\begin{question}[points=5]
$\triangle ABC$中, $D$为$AB$边中点, $\overrightarrow{CE} = \dfrac{1}{3}\overrightarrow{CD}$, $\overrightarrow{AB} = \vec{a}$, $\overrightarrow{AC} = \vec{b}$, 则$\overrightarrow{AE} = $（用$\vec{a}$, $\vec{b}$表示）, 若$|\overrightarrow{AE}| = 5$, $\overrightarrow{AE} \perp \overrightarrow{CB}$, 则$\overrightarrow{AE} \cdot \overrightarrow{CD} = $.
\fillin[{$\dfrac{1}{6}\vec{a} + \dfrac{2}{3}\vec{b}$；$-15$}]
\begin{solution}
\textbf{答案：}$\dfrac{1}{6}\vec{a} + \dfrac{2}{3}\vec{b}$；$-15$

\textbf{分析：}根据向量的线性运算求解即可空一，应用数量积运算律计算求解空二.

\textbf{详解：}如图，因为$\overrightarrow{CE} = \dfrac{1}{3}\overrightarrow{CD}$, 所以$\overrightarrow{AE} - \overrightarrow{AC} = \dfrac{1}{3}(\overrightarrow{AD} - \overrightarrow{AC})$, 所以$\overrightarrow{AE} = \dfrac{1}{3}\overrightarrow{AD} + \dfrac{2}{3}\overrightarrow{AC}$. 因为$D$为线段$AB$的中点, 所以$\overrightarrow{AE} = \dfrac{1}{6}\overrightarrow{AB} + \dfrac{2}{3}\overrightarrow{AC} = \dfrac{1}{6}\vec{a} + \dfrac{2}{3}\vec{b}$. 又因为$|\overrightarrow{AE}| = 5$, $\overrightarrow{AE} \perp \overrightarrow{CB}$, 所以$|\overrightarrow{AE}|^2 = \left(\dfrac{1}{6}\vec{a} + \dfrac{2}{3}\vec{b}\right)^2 = \dfrac{1}{36}|\vec{a}|^2 + \dfrac{2}{9}\vec{a}\cdot\vec{b} + \dfrac{4}{9}|\vec{b}|^2 = 25$, $\overrightarrow{AE} \cdot \overrightarrow{CB} = \left(\dfrac{1}{6}\vec{a} + \dfrac{2}{3}\vec{b}\right)\cdot(\vec{a} - \vec{b}) = \dfrac{1}{6}|\vec{a}|^2 + \dfrac{1}{2}\vec{a}\cdot\vec{b} - \dfrac{2}{3}|\vec{b}|^2 = 0$, 所以$|\vec{a}|^2 + 3\vec{a}\cdot\vec{b} = 4|\vec{b}|^2$. 所以$|\vec{a}|^2 + 4\vec{a}\cdot\vec{b} = 180$, 所以$\overrightarrow{AE} \cdot \overrightarrow{CD} = \left(\dfrac{1}{6}\vec{a} + \dfrac{2}{3}\vec{b}\right)\cdot\left(-\vec{b} + \dfrac{1}{2}\vec{a}\right) = \dfrac{1}{12}|\vec{a}|^2 + \dfrac{1}{6}\vec{a}\cdot\vec{b} - \dfrac{2}{3}|\vec{b}|^2 = \dfrac{1}{12}(|\vec{a}|^2 + 2\vec{a}\cdot\vec{b} - 8|\vec{b}|^2) = \dfrac{1}{12}(|\vec{a}|^2 + 2\vec{a}\cdot\vec{b} - 2(|\vec{a}|^2 + 3\vec{a}\cdot\vec{b})) = \dfrac{1}{12}(-|\vec{a}|^2 - 4\vec{a}\cdot\vec{b}) = -\dfrac{1}{12}(|\vec{a}|^2 + 4\vec{a}\cdot\vec{b}) = -\dfrac{180}{12} = -15$. 故答案为$\dfrac{1}{6}\vec{a} + \dfrac{2}{3}\vec{b}$；$-15$.
\end{solution}
\end{question}

\begin{question}[points=5]
若$a, b \in \mathbb{R}$, 对$\forall x \in [-2, 2]$, 均有$(2a+b)x^2 + bx - a - 1 \le 0$恒成立, 则$2a+b$的最小值为.
\fillin[{$-4$}]
\begin{solution}
\textbf{答案：}$-4$

\textbf{分析：}先设$t = 2a + b$, 根据不等式的形式, 为了消$a$可以取$x = -\dfrac{1}{2}$, 得到$t \ge -4$, 验证$t = -4$时, $a, b$是否可以取到, 进而判断该最小值是否可取即可得到答案.

\textbf{详解：}设$t = 2a + b$, 原题转化为求$t$的最小值. 原不等式可化为对任意的$-2 \le x \le 2$, $tx^2 + (t - 2a)x - a - 1 \le 0$. 不妨代入$x = -\dfrac{1}{2}$, 得$\dfrac{1}{4}t - \dfrac{1}{2}(t - 2a) - a - 1 \le 0$, 得$t \ge -4$. 当$t = -4$时, 原不等式可化为$-4x^2 + (-4 - 2a)x - a - 1 \le 0$, 即$-\left[2x + \left(\dfrac{1}{2}a + 1\right)\right]^2 + \dfrac{1}{4}a^2 \le 0$. 观察可知, 当$a = 0$时, $-(2x+1)^2 \le 0$对$-2 \le x \le 2$一定成立, 当且仅当$x = -\dfrac{1}{2}$取等号. 此时, $a=0$, $b=-4$, 说明$t=-4$时, $a,b$均可取到, 满足题意. 故$t = 2a + b$的最小值为$-4$. 故答案为$-4$.
\end{solution}
\end{question}

\section*{三、解答题}
本大题共5小题，共75分，解答应写出文字说明，证明过程或演算步骤.

\begin{question}[points=15]
在$\triangle ABC$中, 角$A, B, C$的对边分别为$a, b, c$. 已知$a\sin B = \sqrt{3}b\cos A$, $c - 2b = 1$, $a = \sqrt{7}$.

(1) 求$A$的值;

(2) 求$c$的值;

(3) 求$\sin(A + 2B)$的值.
\begin{solution}

\textbf{分析：}(1)由正弦定理化边为角再化简可求；(2)由余弦定理，结合(1)结论与已知代入可得关于$b$的方程，求解可得$b$，进而求得$c$；(3)利用正弦定理先求$B$，再由二倍角公式分别求$\sin 2B$, $\cos 2B$，由两角和的正弦可得.

\textbf{详解：}（1）已知$a\sin B = \sqrt{3}b\cos A$, 由正弦定理$\dfrac{a}{\sin A} = \dfrac{b}{\sin B}$, 得$a\sin B = b\sin A = \sqrt{3}b\cos A$, 显然$\cos A \neq 0$, 得$\tan A = \sqrt{3}$, 由$0 < A < \pi$, 故$A = \dfrac{\pi}{3}$.
（2）由（1）知$\cos A = \dfrac{1}{2}$, 且$c = 2b + 1$, $a = \sqrt{7}$, 由余弦定理$a^2 = b^2 + c^2 - 2bc\cos A$, 则$7 = b^2 + (2b+1)^2 - 2 \times \dfrac{1}{2} b(2b+1) = 3b^2 + 3b + 1$, 解得$b = 1$（$b = -2$舍去）, 故$c = 3$.
（3）由正弦定理$\dfrac{a}{\sin A} = \dfrac{b}{\sin B}$, 且$b = 1$, $a = \sqrt{7}$, $\sin A = \dfrac{\sqrt{3}}{2}$, 得$\sin B = \dfrac{b\sin A}{a} = \dfrac{\sqrt{21}}{14}$, 且$a > b$, 则$B$为锐角, 故$\cos B = \dfrac{5\sqrt{7}}{14}$, 故$\sin 2B = 2\sin B\cos B = \dfrac{5\sqrt{3}}{14}$, 且$\cos 2B = 1 - 2\sin^2 B = 1 - 2 \times \left(\dfrac{\sqrt{21}}{14}\right)^2 = \dfrac{11}{14}$; 故$\sin(A + 2B) = \sin A \cos 2B + \cos A \sin 2B = \dfrac{\sqrt{3}}{2} \times \dfrac{11}{14} + \dfrac{1}{2} \times \dfrac{5\sqrt{3}}{14} = \dfrac{4\sqrt{3}}{7}$.
\end{solution}
\end{question}

\begin{question}[points=15]
正方体$ABCD - A_1B_1C_1D_1$的棱长为4, $E$、$F$分别为$A_1D_1$, $C_1B_1$中点, $CG = 3GC_1$.

(1) 求证: $GF \perp$平面$FBE$;

(2) 求平面$FBE$与平面$EBG$夹角的余弦值;

(3) 求三棱锥$D - FBE$的体积.
\begin{solution}

\textbf{分析：}(1)法一、利用正方形的性质先证明$FG \perp BF$, 再结合正方体的性质得出$EF \perp$平面$BCC_1B_1$, 利用线面垂直的性质与判定定理证明即可；法二、建立空间直角坐标系, 利用空间向量证明线面垂直即可；(2)利用空间向量计算面面夹角即可；(3)利用空间向量计算点面距离，再利用锥体的体积公式计算即可.

\textbf{详解：}（1）法一、在正方形$BCC_1B_1$中，由条件易知$\tan\angle C_1FG = \dfrac{C_1G}{C_1F} = \dfrac{1}{2} = \dfrac{FB_1}{BB_1} = \tan\angle B_1BF$，所以$\angle C_1FG = \angle B_1BF$，则$\angle B_1FB + \angle B_1BF = \dfrac{\pi}{2} = \angle C_1FG + \angle B_1FB$，故$\angle BFG = \pi - (\angle C_1FG + \angle B_1FB) = \dfrac{\pi}{2}$，即$FG \perp BF$. 在正方体中，易知$D_1C_1 \perp$平面$BCC_1B_1$，且$EF // D_1C_1$，所以$EF \perp$平面$BCC_1B_1$，又$FG \subset$平面$BCC_1B_1$，∴$EF \perp FG$. ∵$EF \cap BF = F$, $EF, BF \subset$平面$BEF$，∴$GF \perp$平面$BEF$.
法二、如图以$D$为中心建立空间直角坐标系，则$B(4,4,0)$, $E(2,0,4)$, $F(2,4,4)$, $G(0,4,3)$, 所以$\overrightarrow{EF} = (0,4,0)$, $\overrightarrow{EB} = (2,4,-4)$, $\overrightarrow{FG} = (-2,0,-1)$. 设$\vec{m} = (a,b,c)$是平面$BEF$的一个法向量，则$\begin{cases} \vec{m} \cdot \overrightarrow{EF} = 4b = 0 \\ \vec{m} \cdot \overrightarrow{EB} = 2a + 4b - 4c = 0 \end{cases}$, 令$a=2$, 则$b=0$, $c=1$, 所以$\vec{m} = (2,0,1)$. 易知$\overrightarrow{FG} = -\vec{m}$, 则$\overrightarrow{FG}$也是平面$BEF$的一个法向量，∴$GF \perp$平面$BEF$.
（2）同上法二建立的空间直角坐标系，所以$\overrightarrow{EG} = (-2,4,-1)$, $\overrightarrow{BG} = (-4,0,3)$. 由(1)知$\overrightarrow{FG}$是平面$BEF$的一个法向量. 设平面$BEG$的一个法向量为$\vec{n} = (x,y,z)$, 所以$\begin{cases} \vec{n} \cdot \overrightarrow{EG} = -2x + 4y - z = 0 \\ \vec{n} \cdot \overrightarrow{BG} = -4x + 3z = 0 \end{cases}$, 令$x=6$, 则$z=8$, $y=5$, 即$\vec{n} = (6,5,8)$. 设平面$BEF$与平面$BEG$的夹角为$\alpha$, 则$\cos\alpha = |\cos\langle\overrightarrow{FG},\vec{n}\rangle| = \dfrac{|\overrightarrow{FG}\cdot\vec{n}|}{|\overrightarrow{FG}||\vec{n}|} = \dfrac{20}{\sqrt{5} \times 5\sqrt{5}} = \dfrac{4}{5}$.
（3）由(1)知$EF \perp$平面$BCC_1B_1$, $FB \subset$平面$BCC_1B_1$, ∴$EF \perp FB$. 易知$S_{\triangle BEF} = \dfrac{1}{2} EF \cdot BF = \dfrac{1}{2} \times 4 \times \sqrt{4^2 + 2^2} = 4\sqrt{5}$. 又$\overrightarrow{DE} = (2,0,4)$, 则$D$到平面$BEF$的距离为$d = \dfrac{|\overrightarrow{DE} \cdot \overrightarrow{FG}|}{|\overrightarrow{FG}|} = \dfrac{8}{\sqrt{5}}$. 由棱锥的体积公式知：$V_{D-BEF} = \dfrac{1}{3} d \times S_{\triangle BEF} = \dfrac{1}{3} \times \dfrac{8}{\sqrt{5}} \times 4\sqrt{5} = \dfrac{32}{3}$.
\end{solution}
\end{question}

\begin{question}[points=15]
已知椭圆$\dfrac{x^2}{a^2} + \dfrac{y^2}{b^2} = 1$ ($a > b > 0$)的左焦点为$F$, 右顶点为$A$, $P$为$x = a$上一点, 且直线$PF$的斜率为$\dfrac{1}{3}$, $\triangle PFA$的面积为$\dfrac{3}{2}$, 离心率为$\dfrac{1}{2}$.

(1) 求椭圆的方程;

(2) 过点$P$的直线与椭圆有唯一交点$B$（异于点$A$）, 求证: $PF$平分$\angle AFB$.
\begin{solution}

\textbf{分析：}(1)根据题意，利用椭圆的离心率得到$a=2c$，再由直线$PF$的斜率得到$m=c$，从而利用三角形的面积公式得到关于$c$的方程，解之即可得解；(2)联立直线与椭圆方程，利用其位置关系求得$k$，进而得到直线$PB$的方程与点$B$的坐标，法一：利用向量的夹角公式即可得证；法二：利用两直线的夹角公式即可得证；法三利用正切的倍角公式即可得证；法四：利用角平分线的性质与点线距离公式即可得证.

\textbf{详解：}（1）依题意，设椭圆$\dfrac{x^2}{a^2} + \dfrac{y^2}{b^2} = 1$ ($a > b > 0$)的半焦距为$c$, 则左焦点$F(-c,0)$, 右顶点$A(a,0)$, 离心率$e = \dfrac{c}{a} = \dfrac{1}{2}$, 即$a = 2c$. 因为$P$为$x = a$上一点, 设$P(a,m)$, 又直线$PF$的斜率为$\dfrac{1}{3}$, 则$\dfrac{m - 0}{a - (-c)} = \dfrac{1}{3}$, 即$\dfrac{m}{2c + c} = \dfrac{1}{3}$, 解得$m = c$, 则$P(a,c)$, 即$P(2c,c)$. 因为$\triangle PFA$的面积为$\dfrac{3}{2}$, $|AF| = a - (-c) = a + c = 3c$, 高为$|m| = c$, 所以$S_{\triangle PFA} = \dfrac{1}{2} |AF| |m| = \dfrac{1}{2} \times 3c \times c = \dfrac{3}{2}$, 解得$c = 1$, 则$a = 2c = 2$, $b^2 = a^2 - c^2 = 3$, 所以椭圆的方程为$\dfrac{x^2}{4} + \dfrac{y^2}{3} = 1$.
（2）由(1)可知$P(2,1)$, $F(-1,0)$, $A(2,0)$. 易知直线$PB$的斜率存在, 设其方程为$y = kx + m$, 则$1 = 2k + m$, 即$m = 1 - 2k$. 联立$\begin{cases} y = kx + m \\ \dfrac{x^2}{4} + \dfrac{y^2}{3} = 1 \end{cases}$, 消去$y$得, $(3 + 4k^2)x^2 + 8kmx + 4m^2 - 12 = 0$. 因为直线与椭圆有唯一交点, 所以$\Delta = (8km)^2 - 4(3+4k^2)(4m^2 - 12) = 0$, 即$4k^2 - m^2 + 3 = 0$, 则$4k^2 - (1 - 2k)^2 + 3 = 0$, 解得$k = -\dfrac{1}{2}$, 则$m = 2$, 所以直线$PB$的方程为$y = -\dfrac{1}{2}x + 2$. 联立$\begin{cases} y = -\dfrac{1}{2}x + 2 \\ \dfrac{x^2}{4} + \dfrac{y^2}{3} = 1 \end{cases}$, 解得$\begin{cases} x = 1 \\ y = \dfrac{3}{2} \end{cases}$, 则$B(1, \dfrac{3}{2})$. 以下分别用四种方法证明结论：
法一：则$\overrightarrow{FB} = \left(2, \dfrac{3}{2}\right)$, $\overrightarrow{FP} = (3,1)$, $\overrightarrow{FA} = (3,0)$. 所以$\cos\angle BFP = \dfrac{\overrightarrow{FB} \cdot \overrightarrow{FP}}{|\overrightarrow{FB}||\overrightarrow{FP}|} = \dfrac{2 \times 3 + \dfrac{3}{2} \times 1}{\sqrt{2^2 + \left(\dfrac{3}{2}\right)^2} \times \sqrt{3^2 + 1^2}} = \dfrac{3\sqrt{10}}{10}$, $\cos\angle PFA = \dfrac{\overrightarrow{FA} \cdot \overrightarrow{FP}}{|\overrightarrow{FA}||\overrightarrow{FP}|} = \dfrac{3 \times 3 + 1 \times 0}{3 \times \sqrt{3^2 + 1^2}} = \dfrac{3\sqrt{10}}{10}$, 则$\cos\angle BFP = \cos\angle PFA$, 又$\angle BFP, \angle PFA \in (0, \dfrac{\pi}{2})$, 所以$\angle BFP = \angle PFA$, 即$PF$平分$\angle AFB$.
法二：所以$k_{FB} = \dfrac{\dfrac{3}{2} - 0}{1 - (-1)} = \dfrac{3}{4}$, $k_{PF} = \dfrac{1 - 0}{2 - (-1)} = \dfrac{1}{3}$, $k_{AF} = 0$. 由两直线夹角公式，得$\tan\angle BFP = \left|\dfrac{\dfrac{3}{4} - \dfrac{1}{3}}{1 + \dfrac{1}{3} \times \dfrac{3}{4}}\right| = \dfrac{1}{3}$, $\tan\angle PFA = \left|\dfrac{\dfrac{1}{3} - 0}{1 + 0}\right| = \dfrac{1}{3}$, 则$\tan\angle BFP = \tan\angle PFA$, 又$\angle BFP, \angle PFA \in (0, \dfrac{\pi}{2})$, 所以$\angle BFP = \angle PFA$, 即$PF$平分$\angle AFB$.
法三：则$\tan\angle PFA = k_{PF} = \dfrac{1}{3}$, $\tan\angle BFP = k_{FB} = \dfrac{3}{4}$, 故$\tan 2\angle PFA = \dfrac{2\tan\angle PFA}{1 - \tan^2\angle PFA} = \dfrac{2 \times \dfrac{1}{3}}{1 - \left(\dfrac{1}{3}\right)^2} = \dfrac{3}{4} = \tan\angle BFP$, 又$\angle BFP, \angle PFA \in (0, \dfrac{\pi}{2})$, 所以$\angle BFP = \angle PFA$, 即$PF$平分$\angle AFB$.
法四：则$k_{FB} = \dfrac{\dfrac{3}{2} - 0}{1 - (-1)} = \dfrac{3}{4}$, 所以直线$FB$的方程为$y = \dfrac{3}{4}(x+1)$, 即$3x - 4y + 3 = 0$. 则点$P$到直线$FB$的距离为$d = \dfrac{|3 \times 2 - 4 \times 1 + 3|}{\sqrt{3^2 + (-4)^2}} = 1$, 又点$P$到直线$FA$的距离也为1, 所以$PF$平分$\angle AFB$.
\end{solution}
\end{question}

\begin{question}[points=15]
已知数列$\{a_n\}$是等差数列, $\{b_n\}$是等比数列, $a_1 = b_1 = 2$, $a_2 = b_2 + 1$, $a_3 = b_3$.

(1) 求$\{a_n\}$, $\{b_n\}$的通项公式;

(2) $\forall n \in \mathbb{N}^*$, $I \in \{0,1\}$, 有$T_n = \{p_1 a_1 b_1 + p_2 a_2 b_2 + \cdots + p_{n-1} a_{n-1} b_{n-1} + p_n a_n b_n \mid p_1, p_2, \ldots, p_{n-1}, p_n \in I\}$.

(i) 求证: 对任意实数$t \in T_n$, 均有$t < a_{n+1} b_{n+1}$;

(ii) 求$T_n$所有元素之和.
\begin{solution}

\textbf{分析：}(1)设数列$\{a_n\}$的公差为$d$, 数列$\{b_n\}$公比为$q$ ($q \neq 0$), 由题设列出关于$d$和$q$的方程求解, 再结合等差和等比数列通项公式即可得解；(2)(i)由题意结合(1)求出$a_{n+1}b_{n+1}$和$p_1 a_1 b_1 + p_2 a_2 b_2 + \cdots + p_{n-1} a_{n-1} b_{n-1} + p_n a_n b_n$的最大值, 再作差比较两者大小即可证明；(ii)法一：根据$p_1, p_2, \ldots, p_{n-1}, p_n$中全为1、一个为0其余为1、2个为0其余为1、…、全为0几个情况将$T_n$中的所有元素分系列, 并求出各系列中元素的和, 最后将所有系列所得的和加起来即可得解；法二：根据$T_n$元素的特征得到$T_n$中的所有元素的和中各项$a_i b_i$ ($i \in \{1,2,\ldots,n\}$)出现的次数均为$2^{n-1}$次即可求解.

\textbf{详解：}（1）设数列$\{a_n\}$的公差为$d$, 数列$\{b_n\}$公比为$q$ ($q \neq 0$), 则由题得$\begin{cases} 2 + d = 2q + 1 \\ 2 + 2d = 2q^2 \end{cases}$, 解得$\begin{cases} d = 3 \\ q = 2 \end{cases}$, 所以$a_n = 2 + 3(n-1) = 3n - 1$, $b_n = 2 \times 2^{n-1} = 2^n$.
（2）(i)证明：由(1) $p_n a_n b_n = (3n-1)2^n p_n = 0$或$p_n a_n b_n = (3n-1)2^n > 0$, $a_{n+1}b_{n+1} = (3n+2)2^{n+1}$. 当$p_n a_n b_n = (3n-1)2^n > 0$时, 设$S_n = p_1 a_1 b_1 + p_2 a_2 b_2 + \cdots + p_{n-1} a_{n-1} b_{n-1} + p_n a_n b_n = 2 \times 2 + 5 \times 2^2 + \cdots + (3n-4)2^{n-1} + (3n-1)2^n$, 所以$2S_n = 2 \times 2^2 + 5 \times 2^3 + \cdots + (3n-4)2^n + (3n-1)2^{n+1}$, 所以$-S_n = 4 + 3 \times (2^2 + 2^3 + \cdots + 2^n) - (3n-1)2^{n+1} = 4 + 3 \times \dfrac{2^2(1-2^{n-1})}{1-2} - (3n-1)2^{n+1} = -8 + (4-3n)2^{n+1}$, 所以$S_n = 8 + (3n-4)2^{n+1}$, 为$T_n$中的最大元素, 此时$a_{n+1}b_{n+1} - S_n = (3n+2)2^{n+1} - [8 + (3n-4)2^{n+1}] = 6 \cdot 2^{n+1} - 8 > 0$恒成立, 所以对$\forall t \in T_n$, 均有$t < a_{n+1}b_{n+1}$.
(ii)法一：由(i)得$S_n = 8 + (3n-4)2^{n+1}$, 为$T_n$中的最大元素. 由题意可得$T_n$中的所有元素由以下系列中所有元素组成：当$p_1, p_2, \ldots, p_n$均为1时：此时该系列元素只有$S_n$, 即$\mathrm{C}_n^0$个；当$p_1, p_2, \ldots, p_n$中只有一个为0, 其余均为1时：此时该系列的元素有$S_n - a_1 b_1, S_n - a_2 b_2, \ldots, S_n - a_n b_n$, 共有$\mathrm{C}_n^1$个, 则这$n$个元素的和为$\mathrm{C}_n^1 S_n - (a_1 b_1 + a_2 b_2 + \cdots + a_n b_n) = (\mathrm{C}_n^1 - \mathrm{C}_n^0) S_n$；当$p_1, p_2, \ldots, p_n$中只有2个为0, 其余均为1时：此时该系列的元素为$S_n - a_i b_i - a_j b_j$ ($i,j \in \{1,2,\ldots,n\}$, $i \neq j$), 共有$\mathrm{C}_n^2$个, 则这$n$个元素的和为$\mathrm{C}_n^2 S_n - \mathrm{C}_{n-1}^1 (a_1 b_1 + \cdots + a_n b_n) = (\mathrm{C}_n^2 - \mathrm{C}_{n-1}^1) S_n$；当$p_1, p_2, \ldots, p_n$中有3个为0, 其余均为1时：此时该系列的元素为$S_n - a_i b_i - a_j b_j - a_k b_k$ ($i,j,k \in \{1,2,\ldots,n\}$, $i \neq j \neq k$), 共有$\mathrm{C}_n^3$个, 则这$n$个元素的和为$\mathrm{C}_n^3 S_n - \mathrm{C}_{n-1}^2 (a_1 b_1 + \cdots + a_n b_n) = (\mathrm{C}_n^3 - \mathrm{C}_{n-1}^2) S_n$；依此类推, 当$p_1, p_2, \ldots, p_n$中有$n-1$个为0, 1个为1时：此时该系列的元素为$a_1 b_1, a_2 b_2, \ldots, a_n b_n$, 共有$\mathrm{C}_n^{n-1}$个, 则这$n$个元素的和为$\mathrm{C}_n^{n-1} S_n - \mathrm{C}_{n-1}^{n-2} (a_1 b_1 + \cdots + a_n b_n) = (\mathrm{C}_n^{n-1} - \mathrm{C}_{n-1}^{n-2}) S_n$；当$p_1, p_2, \ldots, p_n$均为0时：此时该系列的元素为0, 即$(\mathrm{C}_n^n - \mathrm{C}_{n-1}^{n-1}) S_n = 0$, 共1个. 综上所述, $T_n$中的所有元素之和为$S_n + (\mathrm{C}_n^1 - 1) S_n + (\mathrm{C}_n^2 - \mathrm{C}_{n-1}^1) S_n + \cdots + (\mathrm{C}_n^{n-1} - \mathrm{C}_{n-1}^{n-2}) S_n + 0 = [(\mathrm{C}_n^0 + \mathrm{C}_n^1 + \cdots + \mathrm{C}_n^n) - (\mathrm{C}_{n-1}^0 + \mathrm{C}_{n-1}^1 + \cdots + \mathrm{C}_{n-1}^{n-1})] S_n = (2^n - 2^{n-1}) S_n = 2^{n-1} S_n = 2^{n-1} \cdot [8 + (3n-4)2^{n+1}]$.
法二：由(i)得$S_n = 8 + (3n-4)2^{n+1}$, 为$T_n$中的最大元素. 由题意可得$T_n = \{S_n, S_n - a_i b_i, S_n - a_i b_i - a_j b_j, \ldots, a_i b_i + a_j b_j, a_i b_i, 0\}$, 所以$T_n$的所有的元素的和中各项$a_i b_i$ ($i \in \{1,2,\ldots,n\}$)出现的次数均为$\mathrm{C}_{n-1}^0 + \mathrm{C}_{n-1}^1 + \cdots + \mathrm{C}_{n-1}^{n-2} + \mathrm{C}_{n-1}^{n-1} = 2^{n-1}$次, 所以$T_n$中的所有元素之和为$2^{n-1} (a_1 b_1 + a_2 b_2 + \cdots + a_n b_n) = 2^{n-1} S_n = 2^{n-1} \cdot [8 + (3n-4)2^{n+1}]$.
\end{solution}
\end{question}

\begin{question}[points=15]
已知函数$f(x) = ax - (\ln x)^2$.

(1) $a = 1$时, 求$f(x)$在点$(1, f(1))$处的切线方程;

(2) $f(x)$有3个零点$x_1, x_2, x_3$且$x_1 < x_2 < x_3$.

(i) 求$a$的取值范围;

(ii) 证明: $(\ln x_2 - \ln x_1) \cdot \ln x_3 < \dfrac{4e}{e-1}$.
\begin{solution}

\textbf{分析：}(1)利用导数的几何意义，求导数值得斜率，由点斜式方程可得；(2)(i)令$f(x)=0$, 分离参数得$a = \dfrac{(\ln x)^2}{x}$, 作出函数$g(x) = \dfrac{(\ln x)^2}{x}$的图象，数形结合可得$a$范围；(ii)由(2)结合图象，可得$x_1, x_2, x_3$范围，整体换元$\ln x_1 = t_1$, $\ln x_2 = t_2$, $\ln x_3 = t_3$, 转化为$\begin{cases} ae^{t_1} = t_1^2 \\ ae^{t_2} = t_2^2 \\ ae^{t_3} = t_3^2 \end{cases}$, 结合②③可得$\begin{cases} \ln a + t_2 = 2\ln t_2 \\ \ln a + t_3 = 2\ln t_3 \end{cases}$, 两式作差, 利用对数平均不等式可得$t_2 t_3 < 4$, 再由$t_1 = ae^{t_1} < a$得$-t_1 < a$, 结合$a = \dfrac{t_3^2}{e^{t_3}}$减元处理, 再构造函数求最值, 放缩法可证明不等式.

\textbf{详解：}（1）当$a=1$时, $f(x) = x - (\ln x)^2$, $x > 0$, 则$f'(x) = 1 - \dfrac{2\ln x}{x}$, 则$f'(1) = 1$, 且$f(1) = 1$, 则切点$(1,1)$, 且切线的斜率为1, 故函数$f(x)$在点$(1, f(1))$处的切线方程为$y = x$.
（2）(i)令$f(x) = ax - (\ln x)^2 = 0$, $x > 0$, 得$a = \dfrac{(\ln x)^2}{x}$. 设$g(x) = \dfrac{(\ln x)^2}{x}$, $x > 0$, 则$g'(x) = \dfrac{2\ln x \cdot x - (\ln x)^2}{x^2} = \dfrac{\ln x (2 - \ln x)}{x^2}$. 由$g'(x) = 0$解得$x=1$或$e^2$, 其中$g(1)=0$, $g(e^2) = \dfrac{4}{e^2}$. 当$0 < x < 1$时, $g'(x) < 0$, $g(x)$在$(0,1)$上单调递减；当$1 < x < e^2$时, $g'(x) > 0$, $g(x)$在$(1, e^2)$上单调递增；当$x > e^2$时, $g'(x) < 0$, $g(x)$在$(e^2, +\infty)$上单调递减；且当$x \to 0$时, $g(x) \to +\infty$；当$x \to +\infty$时, $g(x) \to 0$. 要使函数$f(x)$有3个零点, 则方程$a = g(x)$在$(0, +\infty)$内有3个根, 即直线$y = a$与函数$g(x)$的图象有3个交点. 结合图象可知, $0 < a < \dfrac{4}{e^2}$. 故$a$的取值范围为$(0, \dfrac{4}{e^2})$.
(ii)由图象可知, $0 < x_1 < 1 < x_2 < e^2 < x_3$, 设$\ln x_1 = t_1$, $\ln x_2 = t_2$, $\ln x_3 = t_3$, 则$t_1 < 0 < t_2 < 2 < t_3$, 满足$\begin{cases} ae^{t_1} = t_1^2 & \text{①} \\ ae^{t_2} = t_2^2 & \text{②} \\ ae^{t_3} = t_3^2 & \text{③} \end{cases}$, 由②③可得$\begin{cases} \ln a + t_2 = 2\ln t_2 \\ \ln a + t_3 = 2\ln t_3 \end{cases}$, 两式作差可得$t_3 - t_2 = 2(\ln t_3 - \ln t_2)$, 则由对数均值不等式可得$2 = \dfrac{t_3 - t_2}{\ln t_3 - \ln t_2} > \sqrt{t_2 t_3}$, 则$t_2 t_3 < 4$. 故要证$(\ln x_2 - \ln x_1) \cdot \ln x_3 < \dfrac{4e}{e-1}$, 即证$t_2 t_3 - t_1 t_3 < \dfrac{4e}{e-1}$, 只需证$4 - t_1 t_3 \le \dfrac{4e}{e-1}$, 即证$-t_1 t_3 \le \dfrac{4}{e-1}$. 又因为$t_1 < 0$, $t_1 = ae^{t_1} < a$, 则$-t_1 = |t_1| < a$, 所以$-t_1 t_3 < a t_3 = \dfrac{t_3^2}{e^{t_3}} \cdot t_3 = \dfrac{t_3^2}{e^{t_3}}$, 故只需证$\dfrac{t^2}{e^{t/2}} \le \dfrac{4}{e-1}$ (其中$t = t_3 > 2$). 设函数$\varphi(t) = \dfrac{t^2}{e^{t/2}}$, $t > 2$, 则$\varphi'(t) = \dfrac{2t e^{t/2} - t^2 \cdot \frac{1}{2} e^{t/2}}{e^t} = \dfrac{t(4 - t)}{2e^{t/2}}$. 当$2 < t < 4$时, $\varphi'(t) > 0$, 则$\varphi(t)$在$(2,4)$上单调递增；当$t > 4$时, $\varphi'(t) < 0$, 则$\varphi(t)$在$(4, +\infty)$上单调递减. 故$\varphi(t)_{\max} = \varphi(4) = \dfrac{16}{e^2}$, 即$\varphi(t) \le \dfrac{16}{e^2}$. 而由$4e^2 - 16e + 16 = 4(e-2)^2 > 0$, 可知$\dfrac{16}{e^2} < \dfrac{4}{e-1}$成立, 故命题得证.
\end{solution}
\end{question}

